\section{Recurrence scaling laws：什么时候应该多循环}

能够稳定训练只回答“能不能 loop”，还没有回答“应该 loop 几次”。Parcae 进一步把 recurrence 放进 scaling law：在参数量、数据量和 FLOP budget 之间，寻找 compute-optimal 的循环次数。

\subsection{先回顾参数与数据的经典缩放}

为了理解 recurrence 的位置，本节先回顾经典经验模型，它把 validation loss 写成参数量 $N$ 与数据量 $D$ 的幂律和：

\[
\widehat{L}(N,D)=E+X N^{-\alpha}+Y D^{-\beta}.
\]

在固定训练 FLOPs 下，等 FLOP 曲线的最优点通常随预算向“更大模型、更多数据”共同移动。曲线若向右下方移动，含义不是只增加其中一个，而是两者协同扩大。

\begin{figure}[H]
\centering
\includegraphics[width=0.88\textwidth]{00-55-00-classic-scaling.jpg}
\caption{传统 scaling law：参数量与数据量共同决定 compute-optimal 点。}
\figsource{课堂视频 00:54:13--00:55:44。}
\end{figure}

\subsection{把 recurrence 加入等参数、等 FLOP 曲线}

Parcae 在固定参数量下改变数据量与 recurrence。课堂图中，每条曲线对应一个 FLOP budget；随着数据增多，最优点向更高 recurrence 移动。换句话说，在这些实验范围内，固定模型若训练更多数据，完全保持固定 depth 可能不是最优。

\begin{figure}[H]
\centering
\includegraphics[width=0.88\textwidth]{00-58-30-recurrence-scaling.jpg}
\caption{Parcae 的 recurrence scaling：更多数据对应更高 compute-optimal recurrence。}
\figsource{课堂视频 00:55:44--00:58:31。}
\end{figure}

可以用一个不带具体指数的经验式表达这种关系：

\[
T^{\star}(D;N,C)\approx kD^{\gamma},\qquad \gamma>0,
\]

其中 $T^{\star}$ 是给定参数量 $N$ 与 FLOP budget $C$ 下的最优循环次数，$k,\gamma$ 由实验拟合。讲义故意不填入单个指数，因为它依赖模型规模、数据分布和训练设置，不能从一张课堂图外推到所有模型。

\subsection{同等 FLOPs 下，循环可能优于只加数据}

进一步，课堂把 fixed-depth Transformer 与 compute-optimal looping model 放在相同参数和 FLOP 预算下比较。初步结果显示，把部分额外计算用于 recurrence，而不是只增加数据 token，能获得更低 validation loss 和更好的下游分数。

\begin{figure}[H]
\centering
\includegraphics[width=0.88\textwidth]{00-59-00-loop-vs-fixed-depth.jpg}
\caption{相同参数和 FLOP 预算下，optimal looping 与 fixed-depth 的经验比较。}
\figsource{课堂视频 00:58:31--00:59:30。}
\end{figure}

\begin{warningbox}{不要把初步 scaling law 读成普遍定律}
这些曲线来自有限模型规模和数据设置。它们支持“recurrence 值得成为独立搜索维度”，但还不足以证明所有 frontier pretraining 都应循环。更大的模型、MoE、不同数据 mixture、长上下文和 post-training 都可能改变最优点。
\end{warningbox}

\subsection{本章小结}

Parcae 的 scaling 结果把问题从“loop 或不 loop”改成“在给定参数、数据和 FLOPs 下 loop 多少”。初步证据表明数据与 recurrence 应协同增长，但结论必须在更大规模和更多 workload 上复现。

